% ============================================================================================================
% Abraham's contribution, part A: methods we used that are not in the course textbook (ISLP).
% Place at the end of Section 2 ("Background and Literature Review"), after the Random Forests subsection.
% Each entry says what the method is, why we used it (2-3 sentences), and cites a primary source and, where one
% exists, a textbook. Citation keys are in report/references_abraham.bib; every pointer is checked in
% report/citations.md (where only the paper, not a section, could be verified, the paper is cited as a whole).
% ============================================================================================================

\subsection{Methods Beyond the Course Textbook}
\label{sec:beyond}

The resampling methods described above (cross-validation and the bootstrap) assume that observations are independent. Ours are not: each participant
appears in up to 22 dyads, and every dyad in an event shares the same evening, so outcomes and model errors are
correlated within events. The methods below either respect that structure or regularize models that would otherwise
overfit. For each we say what it is, why we used it, and where it comes from.

\subsubsection{Grouped cross-validation}
In grouped (or blocked) $K$-fold cross-validation, whole groups of observations are assigned to folds, so every
observation of a group is held out together \cite[pp.~917, 919]{Roberts:2017}.
We grouped by event: the 21 events are dealt into five folds, and a model is never tested on an event, or a person, it
was trained on. Random-row folds would let a model recognize participants it met in training and would report the
accuracy of predicting people already seen, not new ones; on our data that inflates the AUC of every model
(Section~\ref{sec:errorbars}).

\subsubsection{Nested cross-validation}
When a model has a tuning parameter, choosing it on the same folds that are used to report performance makes the
reported error optimistic \cite{Varma:2006}. Nested cross-validation selects the parameter by an inner
cross-validation inside each outer training set, then scores the refitted model on the outer held-out fold, so every
choice is made without seeing the test data \cite[Section 7.10.2, pp.~245--247]{Hastie:2009}.
We nest inner folds that are also grouped by event, so tuning never sees a held-out event either.

\subsubsection{Cluster-robust standard errors}
Let $\ell_i$ be a per-dyad quantity, such as a misclassification, and $\bar\ell$ its mean over $n$ dyads in $G$ events. The
cluster-robust (sandwich) variance of $\bar\ell$ sums the residuals within each event before squaring,
\[
\hat V_{\mathrm{CR}}=\frac{G}{G-1}\cdot\frac{1}{n^2}\sum_{g=1}^{G}\Big(\sum_{i\in g}(\ell_i-\bar\ell)\Big)^2,
\]
which allows any correlation within an event and assumes none between events \cite{Liang:1986},
\cite[Section II.C, Eqs.~(10)--(12)]{Cameron:2015}.
The usual standard error treats the 4,184 dyads as independent and is too small by roughly the design effect
$1+(\bar m-1)\rho$, where $\bar m$ is the average cluster size and $\rho$ the within-cluster correlation \cite{Kish:1965};
for a mean this is the Moulton factor of \cite[Section II.B, Eq.~(6)]{Cameron:2015} with $\rho_x=1$.
Because we have only 21 events, we pair it with $t$ quantiles on $G-1$ degrees of freedom and check it against the
bias-reduced CR2 variance \cite[Section 3, Eq.~(7)]{Bell:2002} and the wild cluster bootstrap-$t$, which resamples
Rademacher signs on each event's residual sum \cite{Cameron:2008}; these are the standard few-cluster remedies
\cite[Section VI]{Cameron:2015}. A textbook treatment is \cite[Section 24.5, pp.~829--845]{Cameron:2005}.

\subsubsection{The AUC and its standard error}
The AUC equals the probability that a randomly chosen match receives a higher predicted score than a randomly chosen
non-match \cite[pp.~30--31]{Hanley:1982}, and is estimated by the Mann--Whitney statistic
\[
\hat A=\frac{1}{n_1n_0}\sum_{i:\,y_i=1}\ \sum_{j:\,y_j=0}\Big(\mathbb I(s_i>s_j)+\tfrac12\,\mathbb I(s_i=s_j)\Big).
\]
DeLong, DeLong and Clarke-Pearson estimate its variance from each observation's \emph{placement value}, the share of
the other class it outranks \cite{DeLong:1988}; Obuchowski's version sums the placement values within clusters, so it
allows correlation within an event \cite{Obuchowski:1997}. We use the clustered version for every AUC in this report,
and its paired form for differences between two models scored on the same dyads.

\subsubsection{Repeated cross-validation and the corrected resampled $t$}
Repeating $K$-fold cross-validation $R$ times gives $J=RK$ fold errors $e_1,\dots,e_J$. They are correlated because the
training sets overlap, so the naive variance $S_e^2/J$ is too small, and no unbiased estimator of the variance of
$K$-fold cross-validation exists \cite[Theorem 6, p.~1098]{Bengio:2004}. Nadeau and Bengio's correction inflates it
to $\big(\frac1J+\frac{n_{\mathrm{test}}}{n_{\mathrm{train}}}\big)S_e^2=\big(\frac1J+\frac{1}{K-1}\big)S_e^2$ \cite[Section 3.1, p.~248]{Nadeau:2003}.
We study what this correction estimates when the folds are grouped (Section~\ref{sec:errorbars}).

\subsubsection{Regularized discriminant analysis}
QDA estimates a separate covariance matrix for matches and non-matches. With 74 pre-date variables (124 columns once
categories are one-hot coded) and about 550 matches in a training fold, the match covariance is poorly estimated. Friedman's regularized discriminant analysis
shrinks each class covariance toward a simpler matrix \cite{Friedman:1989}, either the pooled covariance or a
multiple of the identity \cite[Section 4.3.1, Eqs.~(4.13)--(4.14), p.~112]{Hastie:2009}. We use the second kind,
$\widetilde\Sigma_k=(1-r)\hat\Sigma_k+rI$ on standardized predictors, which trades a little bias for a large reduction in
variance.

\subsubsection{The lasso path for logistic regression}
The textbook introduces the lasso for linear regression \cite[Section 6.2.2, Equation 6.7, p.~244]{James:2023}; for
logistic regression the penalty $\lambda\sum_j|\beta_j|$ is added to the negative log-likelihood
\cite[Section 4.4.4, Eq.~(4.31), p.~125]{Hastie:2009}. For the linear lasso, Friedman, Hastie and Tibshirani note that
the smallest penalty that sets every coefficient to zero is $\max_j|\langle x_j,y\rangle|/N$ \cite[Section 2.5, p.~7]{Friedman:2010}.
The same argument for the logistic log-likelihood, whose gradient at the intercept-only fit is $\sum_i x_{ij}(y_i-\bar y)$,
gives $\lambda_{\max}=\max_j\big|\sum_i x_{ij}(y_i-\bar y)\big|$ on the scale of the summed log-likelihood, so we search 30
log-spaced values below it.

\subsubsection{Pre-registration}
Pre-registration means writing down hypotheses, analyses and decision thresholds before seeing the results, so
that confirmatory tests are kept separate from exploration \cite{Nosek:2018}. Every hypothesis in this report was
committed, with its threshold, to our repository's \texttt{PREREG.md} before its test ran, and none was changed
afterwards; tests that failed are reported as failures.

\subsubsection{Simulation studies}
To check an error bar we simulate many data sets from a known truth and record how often the 90\% interval contains
the true value. With $B$ replicates, a coverage near 0.9 has Monte Carlo standard error $\sqrt{0.9\cdot0.1/B}$, about 1.5
percentage points for $B=400$ \cite[Section 5.2, Table 6]{Morris:2019}.
